% ============================================================================
%  BAO CAO BO SUNG (tieng Viet)
%  Vi sao cac phuong phap dung thong ke phu tro suy giam manh hon cac phuong
%  phap khong dung: THUOC DO NGOAI so voi THUOC DO TU THAN.
%
%  File NAY DOC LAP voi VLI_baocao_vi.tex (khong ghi de, khong thay the).
%
%  NGUYEN TAC CUA BAN NAY: khong trich mot con so mo phong nao. Chi gom
%    (a) dai so chung minh duoc, va
%    (b) thuoc tinh cua ham doc tu client/attacks.py, kiem bang phep thu
%        f(lambda*Gamma) =? f(Gamma) -- khong phu thuoc gia dinh ve Q, eps, N.
%
%  Compile:
%     pdflatex VLI_cochesupdo_vi.tex   (chay 3 lan)
%     xelatex  VLI_cochesupdo_vi.tex   (chay 3 lan)
% ============================================================================
\documentclass[11pt,a4paper]{article}

\usepackage{iftex}
\ifPDFTeX
  \usepackage[utf8]{inputenc}
  \usepackage[T5]{fontenc}
\else
  \usepackage{fontspec}
\fi
\usepackage[vietnamese]{babel}

\usepackage[margin=2.3cm]{geometry}
\usepackage{amsmath,amssymb,amsthm,mathtools}
\usepackage{booktabs}
\usepackage{array}
\usepackage[hidelinks]{hyperref}

\newtheorem{dinhly}{Định lý}
\newtheorem{bode}[dinhly]{Bổ đề}
\newtheorem{menhde}[dinhly]{Mệnh đề}
\newtheorem{hequa}[dinhly]{Hệ quả}
\newtheorem{dinhnghia}[dinhly]{Định nghĩa}
\theoremstyle{remark}
\newtheorem{nhanxet}[dinhly]{Nhận xét}
\renewcommand{\proofname}{Chứng minh}

\newcommand{\R}{\mathbb{R}}
\newcommand{\one}{\mathbf{1}}
\newcommand{\pos}[1]{\left[#1\right]_+}
\newcommand{\quant}{\mathcal{Q}}
\newcommand{\epstr}{\bar\varepsilon_{\mathrm{tr}}}
\newcommand{\epsaux}{\bar\varepsilon_{\mathrm{aux}}}
\newcommand{\code}[1]{\texttt{\small #1}}

\renewcommand{\arraystretch}{1.25}

\title{Thước đo ngoài và thước đo tự thân\\[0.3em]
\large Vì sao các phương pháp khôi phục nhãn dùng thống kê phụ trợ\\
suy giảm mạnh hơn các phương pháp không dùng}
\author{}
\date{\today}

\begin{document}
\maketitle

\begin{abstract}
Báo cáo này giải thích một quan sát thực nghiệm: khi mô hình toàn cục được
huấn luyện tốt, các tấn công khôi phục nhãn dùng thống kê phụ trợ (ZLG$^+$,
RLU, LLG$^+$) suy giảm mạnh, còn các tấn công không dùng (LLG, VLI$^*$) thì
không. Luận điểm trung tâm: mọi tấn công loại này đều phải quy một gradient
thành số đếm, nên đều cần một \emph{thang độ} để so sánh. Chúng khác nhau
\emph{chỉ} ở chỗ lấy thang đó từ đâu. Nhóm dùng phụ trợ so tín hiệu với một
thước \emph{cố định} không co theo huấn luyện; nhóm không dùng phụ trợ dựng
thước \emph{từ chính tín hiệu}. Vì độ lớn của gradient lớp cuối \emph{chính
là} sai số huấn luyện (Bổ đề~\ref{bd:dolon}), nhóm thứ nhất thấy tỉ số của
mình co lại, còn nhóm thứ hai không thấy gì thay đổi. Ta chứng minh một
\emph{chặn bẫy} phi tiệm cận (Mệnh đề~\ref{md:bay}): đầu ra của nhóm thứ nhất
bị nhốt trong một quả cầu bán kính $A+C$ quanh một tâm không phụ thuộc nhãn,
đúng ở \emph{mọi} mức sai số --- nên suy giảm là liên tục và bắt đầu sớm, chứ
không phải một ngưỡng. Ta hình thức hóa
điều này bằng \emph{bất biến thang độ}, chứng minh LLG và VLI$^*$ có tính đó
còn ba phương pháp kia không, và kiểm lại bằng một phép thử trên mã nguồn
thật. Cuối cùng ta chứng minh thông tin nhãn nằm trong \emph{hướng} của tín
hiệu chứ không trong độ lớn, nên độ lớn là \emph{tham số nhiễu} --- và việc
so nó với một hằng số ngoài là một lỗi thiết kế, không phải một giới hạn
thông tin.
\end{abstract}

\tableofcontents

% ===========================================================================
\section{Phạm vi và quy tắc trích dẫn}

\subsection{Những gì bản này cố ý không có}

Bản trước của báo cáo này chứa nhiều bảng số sinh từ một mô phỏng tự viết
($Q$ đều ngoài đường chéo, $\bar\varepsilon$ là một vô hướng tự do, $N$ tự
chọn). Những bảng đó không phải bằng chứng: chúng chỉ kiểm tính nhất quán của
đại số với \emph{giả định của chính mô phỏng}. Chúng đã bị loại bỏ hoàn toàn.

Bản này chỉ gồm hai loại khẳng định:

\begin{center}
\begin{tabular}{@{}ll@{}}
\toprule
\textsc{[cm]} & Chứng minh đầy đủ từ định nghĩa, hoặc đọc trực tiếp từ mã nguồn. \\
\textsc{[kt]} & Kiểm bằng \emph{phép thử thuộc tính hàm} trên mã nguồn thật: \\
              & cho cùng đầu vào nhân với $\lambda$ khác nhau, so đầu ra. \\
\bottomrule
\end{tabular}
\end{center}

Phép thử \textsc{[kt]} khác về bản chất với mô phỏng: nó không cần biết
$\Gamma$ sinh ra thế nào, vì nó so $f(\lambda\Gamma)$ với $f(\Gamma)$ trên
cùng một $\Gamma$. Kết quả là một sự thật về \emph{hàm}, đúng với mọi đầu vào.

Mọi đại lượng chưa đo trên pipeline thật được ghi rõ là \textbf{chưa đo}, và
không có con số nào thay cho nó.

\subsection{Câu hỏi}

\begin{quote}
Vì sao ZLG$^+$, RLU, LLG$^+$ suy giảm mạnh khi mô hình được huấn luyện tốt,
trong khi LLG và VLI$^*$ thì không?
\end{quote}

% ===========================================================================
\section{Ký hiệu}
\label{sec:kihieu}

Máy chủ trung thực-nhưng-tò mò quan sát cập nhật của một client sau cửa sổ
gồm $m$ bước SGD cục bộ, mỗi bước trên một lô $B$ mẫu; tổng số nhãn
$n=mB$. Cần khôi phục $N=(N_1,\dots,N_C)$ với $\sum_j N_j=n$.

\begin{center}
\begin{tabular}{@{}cl@{}}
\toprule
\textbf{Ký hiệu} & \textbf{Nghĩa} \\
\midrule
$C$, $B$, $m$, $n=mB$ & số lớp; cỡ lô; số bước cục bộ; tổng số nhãn \\
$y_i$, $e_{y_i}$ & nhãn thật của mẫu $i$; vectơ one-hot của nó \\
$p_i$ & phân phối softmax mô hình gán cho mẫu $i$ \\
$\varepsilon_i = 1-p_{i,y_i}$ & \emph{sai số} của mẫu $i$; $\varepsilon_i\in[0,1)$ \\
$N_j$, $\bar\varepsilon_j$ & số mẫu lớp $j$; sai số trung bình của lớp $j$ \\
$a_j = N_j\bar\varepsilon_j$ & \emph{khối lượng lỗi} của lớp $j$ \\
$\epstr = n^{-1}\sum_i\varepsilon_i$ & sai số huấn luyện trung bình trên cửa sổ \\
$\epsaux$ & sai số trung bình trên tập phụ trợ của kẻ tấn công \\
$\bar p_j = 1-\bar\varepsilon_{\mathrm{aux},j}$ & độ tự tin trên \emph{lớp đúng} $j$, đo trên phụ trợ \\
$s_i = \sum_d e_{i,d}$, $\bar O$ & tổng đặc trưng của mẫu $i$ (sau ReLU, $\ge 0$); trung bình của $s_i$ trên phụ trợ \\
$\Gamma$ & tín hiệu gradient mà một phương pháp đọc được \\
$\quant$ & phép làm tròn về số nguyên \\
\bottomrule
\end{tabular}
\end{center}

Hai dạng tín hiệu xuất hiện trong mã nguồn:
\begin{align}
G &= \sum_{i=1}^{n}\bigl(p_i-e_{y_i}\bigr) && \text{(gradient bias; RLU, LLG, VLI$^*$)},
\label{eq:G}\\
\tilde G_j &= \sum_{i=1}^{n} w_i\bigl(p_{i,j}-\delta_{j,y_i}\bigr),
 \quad w_i=\frac{s_i}{\bar O} && \text{(gradient trọng số; ZLG$^+$, LLG$^+$)}.
\label{eq:Gtilde}
\end{align}

Ma trận nhầm lẫn chuẩn hóa, với $Q_{jj}:=0$:
\begin{equation}
Q_{jk} \;=\; \frac{1}{a_j}\sum_{i:\,y_i=j} p_{i,k}\qquad (k\neq j).
\label{eq:Q}
\end{equation}
Phép chia cho $a_j$ khử $\bar\varepsilon$ khỏi $Q$: tử và mẫu cùng thứ nguyên
$[\varepsilon]^1$. Nên $Q$ mô tả \emph{hình dạng} của lỗi, không mô tả
\emph{độ lớn}. Sự phân biệt này là bản lề của toàn bộ báo cáo.

% ===========================================================================
\section{Hai sự thật về tín hiệu}
\label{sec:haisuthat}

\subsection{Độ lớn của tín hiệu \emph{chính là} sai số}

\begin{bode}[Ma trận $Q$ ngẫu nhiên theo hàng]
\label{bd:Qstoch}
\textsc{[cm]}\; Với mọi $j$ sao cho $a_j>0$: $\sum_k Q_{jk}=1$.
\end{bode}

\begin{proof}
$\sum_{k\ne j}\sum_{i:y_i=j}p_{i,k}=\sum_{i:y_i=j}(1-p_{i,j})
 =\sum_{i:y_i=j}\varepsilon_i=a_j$; chia hai vế cho $a_j$.
\end{proof}

\begin{bode}[Đồng nhất thức tín hiệu]
\label{bd:tinhieu}
\textsc{[cm]}\; Với $a=(a_1,\dots,a_C)$,
\begin{equation}
G \;=\; -\bigl(I-Q^{\!\top}\bigr)\,a .
\label{eq:Gid}
\end{equation}
\end{bode}

\begin{proof}
$G_j=\sum_i p_{i,j}-N_j$. Tách theo nhãn thật:
$\sum_{i:y_i=j}p_{i,j}=N_j-a_j$, và $\sum_{i:y_i=k}p_{i,j}=a_kQ_{kj}$ với
$k\ne j$. Dùng $Q_{jj}=0$ để gộp $k=j$ vào tổng mà không đổi giá trị:
$G_j=(N_j-a_j)+\sum_k a_kQ_{kj}-N_j=-\bigl[(I-Q^{\!\top})a\bigr]_j$.
\end{proof}

\begin{bode}[Độ lớn tín hiệu bị chặn bởi sai số huấn luyện]
\label{bd:dolon}
\textsc{[cm]}\;
\begin{equation}
\|G\|_1 \;\le\; 2n\,\epstr,
\qquad
\|\tilde G\|_1 \;\le\; 2n\,w_{\max}\,\epstr .
\label{eq:dolon}
\end{equation}
\end{bode}

\begin{proof}
Từ \eqref{eq:Gid} và Bổ đề \ref{bd:Qstoch}:
$\|G\|_1\le\|a\|_1+\|Q^{\!\top}a\|_1=2\|a\|_1$, và
$\|a\|_1=\sum_j N_j\bar\varepsilon_j=\sum_i\varepsilon_i=n\epstr$. Với
$\tilde G$: mỗi mẫu góp $w_i\|p_i-e_{y_i}\|_1=2w_i\varepsilon_i
\le 2w_{\max}\varepsilon_i$.
\end{proof}

\begin{nhanxet}[Đây là chỗ không có đường thoát]
\label{nx:khongthoat}
Bổ đề \ref{bd:dolon} không phải một hiện tượng phụ của kiến trúc hay dữ liệu.
Gradient của lớp tuyến tính cuối dưới entropy chéo là \emph{tổng các phần dư}
$p_i-e_{y_i}$, và chuẩn $L_1$ của một phần dư bằng đúng $2\varepsilon_i$. Nên
độ lớn của tín hiệu \emph{chính là} sai số, không phải chỉ tương quan với nó.
Huấn luyện mô hình tốt lên \emph{bắt buộc} làm $\|\Gamma\|_1$ co lại. Mọi
tấn công đều phải sống với điều này; câu hỏi duy nhất là nó có \emph{quan
tâm} hay không.

Giả thiết duy nhất ở đây là $w_{\max}=\max_i s_i/\bar O$ bị chặn, và nó chỉ
cần cho $\tilde G$ (tức ZLG$^+$, LLG$^+$). Giá trị $w_{\max}$ thực tế:
\textbf{chưa đo}.
\end{nhanxet}

\subsection{Tín hiệu có tổng bằng không}

\begin{bode}[Tổng triệt tiêu]
\label{bd:tongkhong}
\textsc{[cm]}\; $\sum_{j=1}^{C}\Gamma_j=0$ với $\Gamma=G$ \emph{và} với
$\Gamma=\tilde G$, bất kể trọng số $w_i$.
\end{bode}

\begin{proof}
Với mỗi mẫu $i$: $\sum_j (p_{i,j}-\delta_{j,y_i})=1-1=0$, vì $p_i$ là phân
phối xác suất và $e_{y_i}$ là one-hot. Mọi tổ hợp tuyến tính theo $i$ giữ
tính chất này.
\end{proof}

\begin{hequa}[Ánh xạ nhãn $\to$ tín hiệu suy biến]
\label{hq:hatnhan}
\textsc{[cm]}\; $\one^{\!\top}(I-Q^{\!\top})=0$, nên $I-Q^{\!\top}$ suy biến
và $a\mapsto G$ có hạng $\le C-1$.
\end{hequa}

Hệ quả này sẽ cần đến ở Mục \ref{sec:thongtin}: tín hiệu xác định khối lượng
lỗi chỉ \emph{đến một bậc tự do}, và bậc đó được ghim lại bằng một lớp vắng
mặt.

% ===========================================================================
\section{Vấn đề thang độ}
\label{sec:thangdo}

\subsection{Mọi tấn công đều cần một thước}

Đầu vào là $\Gamma\in\R^C$ (số thực); đầu ra phải là $\hat N\in\mathbb{Z}^C_{\ge0}$
với $\sum_j\hat N_j=n$. Để biến cái trước thành cái sau, thuật toán buộc phải
trả lời câu hỏi \emph{``một đơn vị nhãn tương ứng với bao nhiêu gradient?''}
Đó là một \textbf{thang độ}, và nó phải đến từ đâu đó.

Toàn bộ sự khác biệt giữa năm phương pháp nằm ở câu trả lời đó.

\subsection{Bất biến thang độ}

\begin{dinhnghia}[Bất biến thang độ]
\label{dn:batbien}
Một tấn công $\mathcal{A}:\R^C\to\mathbb{Z}^C$ gọi là \emph{bất biến thang độ}
nếu
\begin{equation}
\mathcal{A}(\lambda\Gamma)\;=\;\mathcal{A}(\Gamma)
\qquad\text{với mọi } \lambda>0 \text{ và mọi } \Gamma .
\label{eq:batbien}
\end{equation}
\end{dinhnghia}

\begin{nhanxet}[Vì sao đúng là khái niệm cần dùng]
\label{nx:visaobatbien}
Theo Bổ đề \ref{bd:dolon}, tác động \emph{trội} của huấn luyện lên $\Gamma$ là
một phép co nhân tính: $\|\Gamma\|_1$ tỉ lệ với $\epstr$. Một tấn công bất
biến thang độ, \emph{theo định nghĩa}, miễn nhiễm với phép co đó. Một tấn
công không bất biến thì không.

Đây là lý do \eqref{eq:batbien} là tiêu chí đúng, chứ không phải một tính
chất ngẫu nhiên dễ thương. Ta sẽ thấy nó chia năm phương pháp thành hai nhóm
đúng khớp với quan sát thực nghiệm.
\end{nhanxet}

% ===========================================================================
\section{Nhóm dùng thống kê phụ trợ}
\label{sec:aux}

\subsection{Dạng chuẩn, đối chiếu mã nguồn}

\begin{menhde}[Dạng chuẩn]
\label{md:dangchuan}
\textsc{[cm]}\; ZLG$^+$ và bước đầu của RLU đều tính
\begin{equation}
\hat N \;=\; \quant\Bigl(R\bigl(c-\beta\Gamma\bigr)\Bigr),
\qquad R(x)=\frac{n\,\pos{x}}{\|\pos{x}\|_1},
\label{eq:dangchuan}
\end{equation}
với $c\in\R^C_{\ge0}$ \emph{không phụ thuộc nhãn của cửa sổ bị tấn công}.
\end{menhde}

\begin{proof}[Chứng minh (đọc mã nguồn)]
\emph{ZLG$^+$} (\code{client/attacks.py}, \code{ZLGp\_attack}). Mã tính
\[
n_j \;=\; mB\Bigl(\bar p_j-\frac{\mathrm{grad}_j}{\bar O}\Bigr),
\qquad
\mathrm{grad}_j=\sum_d\Delta W_{jd}=\frac{\bar O\,\tilde G_j}{n},
\]
nên $n_j=n\bar p_j-\tilde G_j$; rồi cắt âm, chuẩn hóa về tổng $n$, làm tròn.
Vậy $c=n\bar p$, $\beta=1$, $\Gamma=\tilde G$.

\emph{RLU} (\code{client\_utils.py}, \code{estimated\_entropy\_from\_grad}).
Mã đọc nguyên văn là
\begin{center}
\begin{tabular}{@{}ll@{}}
\code{bias = -np.array(bias)}            & $-G/n$ \\
\code{bias[i] = bias[i] + shift[i]}      & $x=c-G/n$ \\
\code{if bias[i] < 0: bias[i] = 0}       & $\pos{x}$ \\
\code{s = np.sum(abs(bias))}             & $\|\pos{x}\|_1$ \\
\code{solution[i] = round(bias[i]/s * B)}& $\quant\bigl(n\pos{x}/\|\pos{x}\|_1\bigr)$ \\
\end{tabular}
\end{center}
với \code{shift = softmax(new\_mu) * rho\_mean}. Vậy
$c=\bar\rho\cdot\mathrm{softmax}(\bar\mu)$, $\beta=1/n$, $\Gamma=G$.
\end{proof}

\subsection{Thước bị ghim ở $\Theta(1)$}

\begin{menhde}[Chuẩn của thước không co theo huấn luyện]
\label{md:ghim}
\textsc{[cm]}\;
\begin{align}
\text{ZLG}^+:&\quad \|c\|_1 \;=\; n\sum_j\bar p_j
  \;=\; nC\,(1-\epsaux) \;\xrightarrow[\ \epsaux\to0\ ]{}\; nC ,
\label{eq:cZLG}\\
\text{RLU}:&\quad \|c\|_1 \;=\; \bar\rho\sum_j\mathrm{softmax}(\bar\mu)_j
  \;=\; \bar\rho\cdot 1 \;=\; \bar\rho .
\label{eq:cRLU}
\end{align}
\end{menhde}

\begin{proof}
\eqref{eq:cZLG}: $\bar p_j=1-\bar\varepsilon_{\mathrm{aux},j}$ theo định nghĩa
trong mã (\code{pj[k] += mean\_probs[k]} tích lũy xác suất trên \emph{chính
lớp} $k$ của mẫu phụ trợ thuộc lớp $k$). \eqref{eq:cRLU}: vectơ softmax có
tổng bằng $1$ theo định nghĩa.
\end{proof}

\begin{nhanxet}[Đây là trả lời cho một phản biện quan trọng]
\label{nx:phanbien}
Phản biện tự nhiên là: \emph{tập phụ trợ cũng được chấm bằng cùng mô hình, nên
$\epsaux$ cũng giảm --- sao thống kê phụ trợ lại không co?}

Mệnh đề \ref{md:ghim} trả lời chính xác. Tiêu chí không phải ``có dùng phụ
trợ hay không'' mà là \emph{đại lượng phụ trợ đó được neo ở đâu}:

\begin{itemize}
\item ZLG$^+$ dùng $\bar p_j=1-\bar\varepsilon_{\mathrm{aux},j}$. Đây là
  \emph{một trừ} lỗi, nên khi $\epsaux$ giảm thì $\bar p_j$ \textbf{tăng về
  $1$}. Thước không co --- nó phình ra tới trần $nC$.
\item RLU dùng $\bar\rho\cdot\mathrm{softmax}(\bar\mu)$. Phụ trợ quyết định
  \emph{hướng} của $c$, nhưng softmax cộng lại bằng $1$ \emph{theo định
  nghĩa}, nên \emph{độ lớn} $\|c\|_1=\bar\rho$ do siêu tham số $(\gamma,m)$
  định đoạt, hoàn toàn không phụ thuộc mô hình.
\item LLG$^+$ thì \emph{khác}: \code{new\_offset} và \code{new\_impact} đến từ
  \code{estimate\_static\_LLG}, và \code{offset[j] += gradients\_for\_prediction[j]}
  --- offset \emph{chính là} gradient tích lũy trên dữ liệu phụ trợ, nên nó
  $\propto\epsaux$ và \textbf{có co}. Xem Mục \ref{sec:llgp}.
\end{itemize}

Nói gọn: ZLG$^+$ và RLU dựng thước từ \emph{xác suất}, mà xác suất bị chuẩn
hóa về tổng $1$ bất kể mô hình giỏi hay dở. Chính phép chuẩn hóa đó ghim
thước lại.
\end{nhanxet}

\subsection{Hệ quả: không bất biến, và sụp về đâu}

\begin{menhde}[Nhóm phụ trợ không bất biến thang độ]
\label{md:khongbatbien}
\textsc{[cm]}\; Với $c\ne 0$, ánh xạ $\Gamma\mapsto\quant(R(c-\beta\Gamma))$
\emph{không} thỏa \eqref{eq:batbien}. Hơn nữa
\begin{equation}
\lim_{\lambda\to0^+}\ \quant\Bigl(R\bigl(c-\beta\lambda\Gamma\bigr)\Bigr)
\;=\; \quant\bigl(R(c)\bigr),
\label{eq:supve}
\end{equation}
một vectơ \emph{không phụ thuộc chút nào} vào $\Gamma$, tức vào nhãn.
\end{menhde}

\begin{proof}
$c-\beta\lambda\Gamma\to c$ khi $\lambda\to0^+$; $R$ liên tục tại $c$ khi
$\|\pos{c}\|_1>0$, và $\quant$ liên tục trừ trên tập các điểm nửa nguyên.
\end{proof}

Kết hợp với Bổ đề \ref{bd:dolon} ($\|\Gamma\|_1=\Theta(\epstr)$) và Mệnh đề
\ref{md:ghim} ($\|c\|_1$ bị ghim), ta có \textbf{toàn bộ định lý sụp đổ gói
trong một dòng}: huấn luyện đóng vai $\lambda\to0$, và đầu ra tiến về
$\quant(R(c))$.

\begin{nhanxet}[Đây mới chỉ là trường hợp giới hạn]
\label{nx:chigioihan}
Phát biểu \eqref{eq:supve} là tiệm cận, nên tự nó \emph{không} giải thích
được vì sao suy giảm đã thấy rõ khi mô hình còn kém. Mục \ref{ssec:bay} thay
nó bằng Mệnh đề \ref{md:bay}, một chặn đúng ở mọi mức sai số; \eqref{eq:supve}
là trường hợp $A\to0$ của chặn đó.
\end{nhanxet}

Vectơ giới hạn đó là gì, thì phụ thuộc phụ trợ: với ZLG$^+$ là
$n\bar p/\sum_j\bar p_j$, với RLU là $n\cdot\mathrm{softmax}(\bar\mu)$. Cả hai
đều là \textbf{tiên nghiệm lớp của tập phụ trợ của kẻ tấn công}. Nói cách
khác, khi mô hình được huấn luyện tốt, các phương pháp này không còn đoán
nhãn của client --- chúng đoán phân bố lớp của chính dữ liệu phụ trợ.

\subsection{Đại lượng điều khiển}

\begin{dinhnghia}[Tỉ số tín hiệu trên thước]
\label{dn:A}
\begin{equation}
A \;:=\; \frac{n\,\beta\,\|\Gamma\|_1}{\|c\|_1}\;.
\label{eq:A}
\end{equation}
\end{dinhnghia}

\begin{hequa}
\label{hq:Aeps}
\textsc{[cm]}\; $A=\Theta(\epstr)$, vì tử số là $\Theta(n\epstr)$ (Bổ đề
\ref{bd:dolon}) còn mẫu số bị ghim (Mệnh đề \ref{md:ghim}). Với ZLG$^+$, mẫu
số còn \emph{tăng} khi $\epsaux$ giảm, nên $A$ bị bóp từ cả hai phía.
\end{hequa}

\begin{nhanxet}[$\bar O$ không triệt tiêu]
\label{nx:Obar}
Trong chứng minh Mệnh đề \ref{md:dangchuan}, $\bar O$ biến mất ở bước cuối.
Đó là \emph{kế toán ký hiệu}: định nghĩa \eqref{eq:Gtilde} của $\tilde G$ đã
chứa $\bar O$ ở mẫu số của $w_i$. Viết theo đúng đại lượng máy chủ
\emph{quan sát được} thì $\bar O$ hiện ra ở $\beta$:
\[
n_j=n\bar p_j-\frac{n}{\bar O}\mathrm{grad}_j
\qquad\Longrightarrow\qquad
\beta=\frac{n}{\bar O},\quad \Gamma=\mathrm{grad},
\quad\text{nên}\quad A\propto\frac{1}{\bar O}.
\]
Vậy $\bar O$ là một thành phần nữa của thước: ước lượng $\bar O$ càng lớn thì
tín hiệu càng bị nén. Hành vi của $\bar O$ khi huấn luyện (tăng? giảm? đứng
yên?) --- \textbf{chưa đo}, và không suy ra được bằng đại số.
\end{nhanxet}

% ===========================================================================
\section{Vì sao tỉ số nhỏ lại giết tín hiệu}
\label{sec:affine}

Mục này làm rõ rằng ở chế độ thông thường, ánh xạ \eqref{eq:dangchuan} là
\emph{affine chính xác}, nên không có gì ẩn: đầu ra là một hằng số cộng một
nhiễu loạn tỉ lệ với $\Gamma$.

\begin{bode}[Phân rã phần dương]
\label{bd:phanra}
\textsc{[cm]}\; Với mọi $x\in\R^C$:
$\ \|\pos{x}\|_1=\sum_j x_j+\sum_{j:x_j<0}|x_j|$.
\end{bode}

\begin{proof}
$[t]_+=t+[-t]_+$ với mọi $t\in\R$; lấy tổng theo $j$.
\end{proof}

\begin{bode}[Không bị kẹp]
\label{bd:khongkep}
\textsc{[cm]}\; Nếu $n\beta\epstr<\min_j c_j$ thì $x=c-\beta\Gamma$ có mọi
tọa độ dương.
\end{bode}

\begin{proof}
Theo Bổ đề \ref{bd:tinhieu}, $\Gamma_j=-a_j+\sum_k Q_{kj}a_k
\le\sum_k a_k=n\epstr$ (dùng $Q_{kj}\le1$). Vậy
$x_j\ge c_j-\beta n\epstr>0$.
\end{proof}

Hai kết quả sau \emph{chỉ} phục vụ dạng affine chính xác và đạo hàm ở Mục
\ref{sec:thongtin}; kết luận chính (Mệnh đề \ref{md:bay}) không cần đến chúng
--- xem Nhận xét \ref{nx:vaitroH2}.

\begin{hequa}[Ngưỡng cụ thể]
\label{hq:nguong}
\textsc{[cm]}\;
\begin{itemize}
\item \textbf{ZLG$^+$}: $c=n\bar p$, $\beta=1$, nên điều kiện là
  $\epstr<\min_j\bar p_j$, tức
  \[
  \boxed{\;\epstr+\max_j\bar\varepsilon_{\mathrm{aux},j}<1\;}
  \]
  Nó hỏng chỉ khi mô hình gần như chưa huấn luyện (cả hai sai số đều gần
  $1-1/C$); thỏa rộng rãi với mọi mô hình đã huấn luyện --- tức đúng trong
  vùng cần giải thích.
\item \textbf{RLU}: $\beta=1/n$ và $\min_j c_j=\bar\rho\min_j\mathrm{softmax}(\bar\mu)_j$,
  nên điều kiện là $\epstr<\bar\rho\min_j\mathrm{softmax}(\bar\mu)_j$. Nếu
  $\bar\mu$ lệch mạnh thì ngưỡng này có thể chặt; giá trị thực tế của
  $\min_j\mathrm{softmax}(\bar\mu)_j$ --- \textbf{chưa đo}.
\end{itemize}
\end{hequa}

\begin{menhde}[Đồng nhất thức affine]
\label{md:affine}
\textsc{[cm]}\; Giả sử \textbf{(H1)} $c\ge0$ và \textbf{(H2)} $x=c-\beta\Gamma$
không có tọa độ âm. Khi đó
\begin{equation}
R\bigl(c-\beta\Gamma\bigr)
\;=\;\underbrace{\frac{n}{\|c\|_1}\,c}_{\text{không phụ thuộc nhãn}}
\;-\;\underbrace{\frac{n\beta}{\|c\|_1}\,\Gamma}_{\text{toàn bộ tín hiệu}}
\label{eq:affine}
\end{equation}
\emph{chính xác}, không phải xấp xỉ.
\end{menhde}

\begin{proof}
Theo (H2), $\pos{x}=x$, nên Bổ đề \ref{bd:phanra} cho
$\|\pos{x}\|_1=\sum_j x_j$. Theo Bổ đề \ref{bd:tongkhong},
$\sum_j x_j=\sum_j c_j$. Theo (H1), $\sum_j c_j=\|c\|_1$. Vậy mẫu số của $R$
là hằng số $\|c\|_1$, và $R(x)=nx/\|c\|_1$.
\end{proof}

\begin{nhanxet}[Hai giả thiết làm hai việc khác nhau]
\label{nx:haigiathiet}
Đây là chỗ dễ sai nhất. $\sum_j\Gamma_j=0$ \textbf{không} kéo theo
$\|\pos{c-\beta\Gamma}\|_1=\|c\|_1$: nó chỉ cho tổng \emph{có dấu}
$\sum_j x_j=\|c\|_1$. Cần thêm (H2) để $\|\pos{x}\|_1$ bằng tổng có dấu đó.
Phân công rõ ràng:
\begin{itemize}
\item \textbf{(H2) không kẹp} làm bước $\|\pos{x}\|_1=\sum_j x_j$ --- thuần
  hình học của $\pos{\cdot}$, không liên quan gradient.
\item \textbf{Bổ đề \ref{bd:tongkhong}} làm bước $\sum_j x_j=\sum_j c_j$ ---
  đây là nơi nhãn và $\epstr$ rời khỏi mẫu số.
\end{itemize}
Không cái nào thay được cái kia. Khi (H2) hỏng, công thức vẫn đóng nhưng
mẫu số lớn hơn: $\|c\|_1+\sum_{j:x_j<0}|x_j|$ theo Bổ đề \ref{bd:phanra}.
\end{nhanxet}

\subsection{Chặn bẫy: phát biểu phi tiệm cận}
\label{ssec:bay}

Mệnh đề \ref{md:khongbatbien} là một phát biểu \emph{giới hạn}
($\lambda\to0^+$). Mục này thay nó bằng một chặn đúng ở \emph{mọi} mức sai
số, không cần ngưỡng --- đây mới là dạng khớp với quan sát thực nghiệm rằng
suy giảm đã xảy ra khi $\epstr$ còn cao.

\begin{menhde}[Đầu ra bị nhốt quanh một tâm không phụ thuộc nhãn]
\label{md:bay}
\textsc{[cm]}\; Chỉ cần \textbf{(H1)} $c\ge0$, $c\ne0$ --- \textbf{không} cần
\textbf{(H2)}. Với $A$ như ở Định nghĩa \ref{dn:A}:
\begin{equation}
\bigl\|R(c-\beta\Gamma)-R(c)\bigr\|_1 \;\le\; A ,
\label{eq:dichchuyen}
\end{equation}
và \emph{đẳng thức xảy ra khi và chỉ khi} \textbf{(H2)} đúng. Do đó, sau khi
làm tròn,
\begin{equation}
\bigl\|\,\hat N-\quant\bigl(R(c)\bigr)\,\bigr\|_1 \;\le\; A+C ,
\label{eq:bay}
\end{equation}
trong đó tâm $\quant(R(c))$ \textbf{không phụ thuộc nhãn} của cửa sổ bị tấn
công.
\end{menhde}

\begin{proof}
Đặt $g=\beta\Gamma$, $L=\|c\|_1$ và $\varphi(t)=R(c-tg)$ với $t\in[0,1]$. Ký
hiệu $S(t)=\{j:c_j-tg_j>0\}$, $g_S$ là $g$ thu hẹp về $S$, và
$D(t)=\|\pos{c-tg}\|_1=\sum_{j\in S(t)}(c_j-tg_j)$.

\emph{Bước 1 ($D\ge L$).} Theo Bổ đề \ref{bd:tongkhong},
$\sum_j(c_j-tg_j)=\sum_j c_j=L$ (dùng \textbf{(H1)}). Theo Bổ đề
\ref{bd:phanra}, $D(t)=L+\sum_{j\notin S(t)}|c_j-tg_j|\ge L>0$. Nói riêng
$\varphi$ xác định trên cả đoạn, không cần \textbf{(H2)}.

\emph{Bước 2 (dấu ngoài $S$).} Nếu $j\notin S(t)$ thì $c_j-tg_j\le0$, mà
$c_j\ge0$, nên $tg_j\ge c_j\ge0$: \emph{mọi tọa độ bị kẹp đều có $g_j\ge0$.}

\emph{Bước 3 (đạo hàm từng khúc).} Mỗi $c_j-tg_j$ affine theo $t$ nên đổi dấu
nhiều nhất một lần; $S(t)$ chỉ thay đổi tại hữu hạn điểm, và $\varphi$ liên
tục, Lipschitz từng khúc. Trên một khúc, đặt $\sigma=\sum_{j\in S}g_j$ thì
$D'=-\sigma$ và
\[
\varphi' \;=\; -\frac{n}{D}\,g_S+\frac{\sigma}{D}\,\varphi,
\qquad\text{nên}\qquad
\|\varphi'\|_1\;\le\;\frac{n}{D}\Bigl(\|g_S\|_1+|\sigma|\Bigr),
\]
dùng $\|\varphi\|_1=n$ (ảnh của $R$ luôn nằm trên đơn hình khối lượng $n$).

\emph{Bước 4 (chỗ then chốt).} Theo Bổ đề \ref{bd:tongkhong},
$\sigma=-\sum_{j\notin S}g_j$; theo Bước 2 mọi số hạng đó không âm, nên
$|\sigma|=\|g_{S^c}\|_1$ --- \emph{bằng}, không phải chặn. Vậy
$\|g_S\|_1+|\sigma|=\|g_S\|_1+\|g_{S^c}\|_1=\|g\|_1$, và cùng Bước 1:
\[
\|\varphi'(t)\|_1\;\le\;\frac{n\|g\|_1}{D(t)}\;\le\;\frac{n\beta\|\Gamma\|_1}{L}\;=\;A .
\]
Lấy tích phân trên $[0,1]$ được \eqref{eq:dichchuyen}. Đẳng thức buộc
$D(t)\equiv L$, tức không tọa độ nào bị kẹp, tức \textbf{(H2)}.

\emph{Phần làm tròn.} $\quant(u)=u+r(u)$ với $\|r(u)\|_\infty\le\tfrac12$, nên
$\|\quant(u)-\quant(v)\|_1\le\|u-v\|_1+C$.
\end{proof}

\begin{nhanxet}[(H2) dùng để làm gì, và không dùng để làm gì]
\label{nx:vaitroH2}
Đây là chỗ dễ hiểu nhầm. Giả thiết không-kẹp \textbf{(H2)} \emph{chỉ} cần cho:
\begin{itemize}
\item \textbf{dạng affine chính xác} \eqref{eq:affine} --- khi có tọa độ bị
  kẹp, mẫu số thành $\|c\|_1+\sum_{j:x_j<0}|x_j|$ (Bổ đề \ref{bd:phanra}),
  phụ thuộc $\Gamma$, nên ánh xạ chỉ còn affine \emph{từng khúc};
\item \textbf{đạo hàm} \eqref{eq:DR} ở Mục \ref{sec:thongtin}, vốn cần một
  khúc duy nhất quanh $c$.
\end{itemize}
Nó \textbf{không} cần cho kết luận --- chặn bẫy \eqref{eq:bay} đúng vô điều
kiện. Hệ quả thực tế: luận điểm của Mục này \emph{không} treo vào Hệ quả
\ref{hq:nguong}, nên việc $\min_j\mathrm{softmax}(\bar\mu)_j$ của RLU chưa đo
không còn là lỗ hổng.

Thêm một điểm dễ hiểu ngược: kẹp \emph{không} phải chế độ hỏng. Tọa độ
$x_j<0$ nghĩa là $\beta\Gamma_j>c_j\ge0$, tức lớp $j$ có gradient dương lớn
--- dấu hiệu của lớp \emph{vắng mặt}; cắt nó về $0$ cho đúng $\hat N_j=0$.
Theo Bước 1, kẹp chỉ làm $D$ tăng, tức \emph{nén} tín hiệu thêm, nên nó chỉ
thu nhỏ quả cầu bẫy chứ không nới.
\end{nhanxet}

\begin{hequa}[Chặn dưới cho sai số khôi phục]
\label{hq:chanduoi}
\textsc{[cm]}\;
\[
\|\hat N-N\|_1\;\ge\;
\underbrace{\bigl\|N-\quant(R(c))\bigr\|_1}_{\text{nhãn thật cách tiên nghiệm phụ trợ bao xa}}
\;-\;(A+C).
\]
Tức là: \textbf{tấn công không thể tốt hơn ``đoán tiên nghiệm lớp của tập phụ
trợ'' quá $A+C$} --- đúng với \emph{mọi} $\epstr$, không cần ngưỡng nào.
\end{hequa}

\begin{proof}
Bất đẳng thức tam giác đảo áp vào \eqref{eq:bay}.
\end{proof}

\begin{nhanxet}[Suy giảm liên tục và cắn sớm, không phải một ngưỡng]
\label{nx:lientuc}
\eqref{eq:bay} nói: tỉ lệ khối lượng đầu ra mà \emph{nhãn} được phép điều
khiển nhiều nhất là $(A+C)/n$. Với ZLG$^+$ ($\beta=1$,
$\|c\|_1=nC(1-\epsaux)$), ghép với Bổ đề \ref{bd:dolon}:
\begin{equation}
\frac{A}{n}\;\le\;\frac{2\,w_{\max}\,\epstr}{C\,(1-\epsaux)} .
\label{eq:tilekhoiluong}
\end{equation}
Chú ý $1/C$ ở mẫu số --- đó là lý do chặn này cắn \emph{sớm}. Thế các giá trị
minh họa $C=10$, $\epsaux=0{,}25$, $w_{\max}=1$, $n=640$ (đều \textbf{chưa
đo} trên pipeline thật, đưa vào chỉ để thấy bậc độ lớn):

\begin{center}
\begin{tabular}{@{}ccc@{}}
\toprule
$\epstr$ & $(A+C)/n$ tối đa & khối lượng do tiên nghiệm phụ trợ quyết định \\
\midrule
$0{,}7$  & $0{,}20$ & $\ge 80\%$ \\
$0{,}3$  & $0{,}10$ & $\ge 90\%$ \\
$0{,}05$ & $0{,}03$ & $\ge 97\%$ \\
\bottomrule
\end{tabular}
\end{center}

Không có ngưỡng và không có sụp đổ đột ngột: suy giảm bắt đầu ngay khi
$\epstr$ còn cao, và chỉ sâu dần. Đây là chỗ báo cáo trả lời phản biện rằng
RLU và ZLG$^+$ đã giảm từ lúc mô hình còn kém.

\eqref{eq:tilekhoiluong} còn cho một \textbf{tiên đoán phân biệt được}: vì
$A/n\propto1/C$, ở cùng $\epstr$ thì CIFAR100 phải suy giảm nặng hơn CIFAR10
khoảng một bậc.
\end{nhanxet}

\begin{nhanxet}[Với RLU chặn này yếu, và ở sai số cao thì rỗng]
\label{nx:rluyeu}
\textsc{[cm]}\; Với RLU, $\beta=1/n$ và $\|c\|_1=\bar\rho$, nên
$A/n\le2\epstr/\bar\rho$ --- \emph{không} có thừa số $1/C$. Và $\bar\rho$
\textbf{không phải đại lượng chưa đo}: \code{rlu\_rho\_mean(args)} chỉ đọc
\code{local\_epochs}, \code{gamma}, \code{lr}, \code{mu}, \code{fl\_scheme};
nó không chạm vào mô hình lẫn tập phụ trợ, nên $\bar\rho$ là \emph{hằng số
qua huấn luyện}. Đây chính là chỗ bịt đường thoát ``độ dốc $n/\|c\|_1$ tăng
lên để bù $\|\Gamma\|_1$ co''. Với $\gamma=0{,}1$, $m=10$ thì
$\bar\rho\approx1{,}0988$.

Hệ quả: độ dốc của hai phương pháp lệch nhau bốn bậc.

\begin{center}
\begin{tabular}{@{}lccc@{}}
\toprule
 & độ dốc $n/\|c\|_1$ & $A/n$ & tại $\epstr=0{,}7$ \\
\midrule
ZLG$^+$ & $1/\bigl(C(1-\epsaux)\bigr)\approx0{,}133$ & $\le0{,}267\,\epstr$ & $0{,}19$ \\
RLU     & $n/\bar\rho\approx583$ & $\le1{,}82\,\epstr$ & $1{,}27$ \\
\bottomrule
\end{tabular}
\end{center}

Ở $\epstr=0{,}7$ chặn của RLU \textbf{rỗng} (quả cầu to hơn cả đơn hình).
Nói thẳng: \emph{Mệnh đề \ref{md:bay} giải thích được ZLG$^+$ suy giảm sớm,
nhưng không đóng được hồ sơ RLU ở vùng sai số cao}; ở đó lời giải thích hiện
chỉ còn dựa vào phát biểu giới hạn \eqref{eq:supve}, yếu hơn hẳn.

Một đường thoát khả dĩ, \textbf{chưa kiểm}: $\|\Gamma\|_1\le2n\epstr$ là chặn
xấu nhất, bỏ qua triệt tiêu giữa $a$ và $Q^{\!\top}a$ trong \eqref{eq:Gid}.
Nếu $\|\Gamma\|_1$ thực tế nhỏ hơn chặn đó vài lần thì RLU cũng rơi vào vùng
chặn cắn. Đo $\|\Gamma\|_1$ thật là ưu tiên cao (Mục \ref{sec:nguyenly}).
\end{nhanxet}

\begin{nhanxet}[Thang độ gây ra lỗi; lượng tử hóa khóa lỗi lại]
\label{nx:luongtu}
Bản trước của mục này quy toàn bộ thất bại cho $\quant$. Đó là \emph{sai địa
chỉ}. Hai sự thật tách rời nhau:

\begin{itemize}
\item \textbf{$A$ nhỏ, một mình nó, đã làm đầu ra sai.} Theo
  \eqref{eq:dichchuyen}, ngay \emph{trước} khi làm tròn, đầu ra đã bị nhốt
  cách $R(c)$ đúng $A$. Nếu $N$ cách tiên nghiệm phụ trợ $\Theta(n)$ --- đúng
  với phân hoạch Dirichlet lệch --- thì tấn công đã sai, kể cả khi tính toán
  bằng số thực chính xác tuyệt đối. $\quant$ \textbf{không cần thiết} cho
  thất bại.
\item \textbf{$\quant$ làm thất bại đó không sửa được.} Phần tuyến tính của
  \eqref{eq:affine} là $-\tfrac{n\beta}{\|c\|_1}I$, \emph{khả nghịch} với mọi
  $\beta>0$: về mặt thông tin, bước số thực không mất gì, dù $A$ nhỏ đến đâu
  --- tín hiệu chỉ bị nén sai thang, và về nguyên tắc còn nhân lại được.
  $\quant$ cắt mất khả năng ấy: khi nhiễu loạn không đẩy nổi một tọa độ qua
  trọn một đơn vị, nó biến mất vĩnh viễn.
\end{itemize}

Nói gọn: \textbf{thang độ sai gây ra lỗi; lượng tử hóa khóa lỗi lại.} Phân
biệt này quan trọng vì nó chỉ đúng chỗ cần sửa --- bỏ $\quant$ đi không cứu
được ZLG$^+$ hay RLU; chỉ bỏ thước ngoài mới cứu được.

Quan hệ định lượng giữa $A$ và irec trên pipeline thật --- \textbf{chưa đo};
nhưng \eqref{eq:bay} là một ràng buộc \emph{cứng}, kiểm trực tiếp được (Mục
\ref{sec:nguyenly}).
\end{nhanxet}

% ===========================================================================
\section{Nhóm không dùng thống kê phụ trợ}
\label{sec:noaux}

Hai phương pháp dưới đây dựng thước \emph{từ chính tín hiệu}. Chứng minh là
đọc mã nguồn.

\subsection{LLG}

\begin{menhde}[LLG bất biến thang độ]
\label{md:llg}
\textsc{[cm]}\; Hàm \code{LLG\_attack} thỏa \eqref{eq:batbien}.
\end{menhde}

\begin{proof}
Mã (\code{client/attacks.py}) thực hiện, với $g:=\Gamma$:
\begin{enumerate}
\item $h_1=\{j:g_j<0\}$ --- \emph{bất biến} dưới $g\mapsto\lambda g$,
  $\lambda>0$ (chỉ phụ thuộc dấu).
\item \code{negative\_gradient} $=\sum_{j\in h_1}g_j$ --- \emph{thuần nhất
  bậc 1}.
\item \code{new\_impact} $=(1+1/C)\cdot$\code{negative\_gradient}$/B$ ---
  thuần nhất bậc 1.
\item Vòng lặp cập nhật $g_j \mathrel{-}= $ \code{new\_impact}: cả $g$ lẫn
  bước nhảy cùng thuần nhất bậc 1, nên vectơ $g$ ở mọi thời điểm là thuần
  nhất bậc 1 theo $\Gamma$ ban đầu.
\item \code{argmin} của một vectơ thuần nhất bậc 1 là \emph{bất biến} với
  $\lambda>0$ (phép nhân số dương bảo toàn thứ tự).
\item Đầu ra là \code{prediction.count(i)} --- chỉ phụ thuộc dãy chỉ số đã
  chọn. Bất biến.
\end{enumerate}
\end{proof}

Điểm quyết định là bước (3): LLG lấy bước nhảy từ \code{negative\_gradient},
tức \emph{từ chính dữ liệu bị tấn công}.

\subsection{VLI$^*$}

\begin{menhde}[VLI$^*$ bất biến thang độ]
\label{md:vlistar}
\textsc{[cm]}\; Hàm \code{VLI\_star\_closed\_form} thỏa \eqref{eq:batbien}.
\end{menhde}

\begin{proof}
Toàn bộ phần thân dùng tín hiệu qua đúng một biểu thức:
\begin{center}
\code{g\_max = max(signal)};\qquad
\code{raw = (n\_samples/C) * (1.0 - signal/g\_max)}
\end{center}
Dưới $\mathrm{signal}\mapsto\lambda\,\mathrm{signal}$ với $\lambda>0$: tử số
và \code{g\_max} cùng nhân $\lambda$, nên tỉ số
$\mathrm{signal}/\code{g\_max}$ \emph{không đổi}; do đó \code{raw} không đổi,
và \code{\_integerize\_counts(raw, n)} cho cùng kết quả.
\end{proof}

\begin{nhanxet}[VLI$^*$ đổi chế độ hỏng, không phải xóa chế độ hỏng]
\label{nx:vlidieukien}
Mã \code{raise ValueError} nếu \code{g\_max} $\le$ \code{eps*scale*16}, và
docstring ghi rõ nó cần ít nhất một lớp vắng mặt để \code{g\_max} mang nghĩa
``lớp vắng''. Thật vậy, với lớp $j$ vắng mặt thì $G_j=\sum_i p_{i,j}>0$.

Vậy VLI$^*$ thoát được chế độ hỏng ``co theo $\epstr$'' nhưng nhận một điều
kiện tiên quyết khác: \emph{phải có lớp vắng mặt}. Đó là một đánh đổi có ý
thức, không phải một sự miễn nhiễm vô điều kiện. Mục \ref{sec:thongtin} cho
công thức chính xác cho trường hợp không có lớp vắng.
\end{nhanxet}

\subsection{LLG\texorpdfstring{$^+$}{+}: cùng thuật toán, khác một dòng}
\label{sec:llgp}

\begin{menhde}[LLG$^+$ không bất biến]
\label{md:llgp}
\textsc{[cm]}\; \code{LLGp\_attack} không thỏa \eqref{eq:batbien}.
\end{menhde}

\begin{proof}
Mã thực hiện
\begin{center}
\code{gradients\_for\_prediction -= new\_offset}
\end{center}
với \code{new\_offset} lấy từ \code{estimate\_static\_LLG}, độc lập với
$\Gamma$. Mà $\lambda\Gamma-\text{offset}\ne\lambda(\Gamma-\text{offset})$ khi
$\text{offset}\ne0$, và \code{argmin} của hai vectơ đó khác nhau. Tương tự với
\code{new\_impact}.
\end{proof}

\begin{nhanxet}[Một thí nghiệm có đối chứng nằm sẵn trong mã nguồn]
\label{nx:doichung}
So sánh hai hàm:
\begin{center}
\begin{tabular}{@{}lll@{}}
\toprule
& \textbf{bước nhảy lấy từ} & \textbf{bất biến?} \\
\midrule
\code{LLG\_attack}  & \code{negative\_gradient} (chính tín hiệu) & \textbf{có} \\
\code{LLGp\_attack} & \code{estimate\_static\_LLG(aux)}          & \textbf{không} \\
\bottomrule
\end{tabular}
\end{center}
Ngoài chỗ đó, hai thuật toán gần như trùng nhau: cùng tập dấu $h_1$, cùng vòng
lặp \code{argmin} tham lam, cùng cách đếm. Nên cặp LLG\,/\,LLG$^+$ cô lập
đúng một biến: \emph{thước lấy từ dữ liệu hay từ phụ trợ}.
\end{nhanxet}

\begin{nhanxet}[LLG$^+$ bị điều khiển bởi một tỉ số]
\label{nx:llgpratio}
\textsc{[cm]}\; Trong \code{estimate\_static\_LLG}, offset được tích lũy bằng
\begin{center}
\code{offset[j] += gradients\_for\_prediction[j]}
\end{center}
--- nó \emph{chính là}
gradient trên dữ liệu phụ trợ, nên offset $\propto\epsaux$ theo Bổ đề
\ref{bd:dolon}. Tín hiệu thì $\propto\epstr$. Vậy tỉ số tín hiệu/thước của
LLG$^+$ là
\[
\frac{\epstr}{\epsaux},
\]
tức \textbf{khoảng cách tổng quát hóa}, không phải $\epstr$. Hệ quả kiểm
chứng được: nếu lỗi train và lỗi phụ trợ co cùng nhịp, LLG$^+$ \emph{không}
sụp theo cơ chế này; nó chỉ suy giảm đúng theo mức mô hình overfit.

Trong mã, tập phụ trợ lấy từ \emph{test set}
(\code{aux\_dataset = LocalDataset(test\_dataset, aux\_dict)} trong
\code{main.py}), nên khoảng cách này thực sự tồn tại. Độ lớn của nó ---
\textbf{chưa đo}.
\end{nhanxet}

% ===========================================================================
\section{Kiểm chứng: phép thử thuộc tính hàm}
\label{sec:kiemchung}

\subsection{Thiết kế phép thử}

Với mỗi hàm tấn công $\mathcal{A}$ lấy từ \code{client/attacks.py} và
\code{client/client\_utils.py}, ta cố định một $\Gamma$ rồi so
$\mathcal{A}(\lambda\Gamma)$ với $\mathcal{A}(\Gamma)$ trên dãy
$\lambda\in\{1,10^{-1},10^{-2},10^{-3},10^{-5},10^{-8}\}$.

Đây \emph{không} phải mô phỏng: phép thử không cần biết $\Gamma$ sinh ra thế
nào, không giả định gì về $Q$, $\epstr$ hay $N$. Nó kiểm một đẳng thức giữa
hai lần gọi cùng một hàm, nên kết quả là tính chất của \emph{hàm}.

\subsection{Kết quả}

Với $C=10$, $B=40$, $m=10$, $n=400$, và $\Gamma$ sinh ngẫu nhiên có tổng $0$
(ba lần thử độc lập): \textsc{[kt]}

\begin{center}
\begin{tabular}{@{}lcc@{}}
\toprule
\textbf{Hàm} & \textbf{thước lấy từ} & \textbf{bất biến trên} \\
\midrule
\code{LLG\_attack}                   & chính tín hiệu & $3/3$ phép thử \\
\code{VLI\_star\_closed\_form}       & chính tín hiệu & $3/3$ phép thử \\
\code{estimated\_entropy\_from\_grad} (lõi RLU) & phụ trợ $+$ $\bar\rho$ & $0/3$ phép thử \\
\bottomrule
\end{tabular}
\end{center}

LLG và VLI$^*$ cho đầu ra \emph{giống nhau từng số nguyên} tại mọi $\lambda$
xuống tới $10^{-8}$. Lõi RLU đổi đầu ra ở mọi bước, và hội tụ về vectơ
\emph{đều} $n/C$ --- tương thích với \eqref{eq:supve} khi
$\mathrm{shift}$ được đặt đều trong phép thử.

LLG$^+$ và ZLG$^+$ không được thử bằng số vì chúng gọi
\code{estimate\_static\_*} cần GPU; tính không bất biến của chúng là đại số
hiển nhiên (Mệnh đề \ref{md:llgp}) nên không cần thử.

\begin{nhanxet}[Phép thử này chứng minh và không chứng minh điều gì]
\label{nx:gioihankiem}
Nó chứng minh: LLG và VLI$^*$ miễn nhiễm với \emph{phép co nhân tính} của tín
hiệu; nhóm phụ trợ thì không, và hội tụ về một vectơ không phụ thuộc nhãn.

Nó \emph{không} chứng minh LLG và VLI$^*$ miễn nhiễm với huấn luyện nói
chung. Huấn luyện không chỉ làm co $\Gamma$ --- nó còn làm \emph{trôi hướng}
của $\Gamma$, vì $Q$ thay đổi. Phép thử giữ hướng cố định và chỉ đổi độ lớn,
nên nó cô lập đúng một hiệu ứng. Bất biến thang độ loại bỏ \textbf{một} chế
độ hỏng --- chế độ trội --- không phải mọi chế độ.
\end{nhanxet}

% ===========================================================================
\section{Thông tin nằm ở hướng, không ở độ lớn}
\label{sec:thongtin}

Mục này cho thấy sụp đổ là \emph{lựa chọn thiết kế}, không phải giới hạn
thông tin --- và do đó biện minh cho nguyên lý ở Mục \ref{sec:nguyenly}.

\begin{menhde}[Hướng được bảo toàn hoàn toàn]
\label{md:huong}
\textsc{[cm]}\; Trên siêu phẳng $H=\{v:\sum_j v_j=0\}$ --- nơi $\Gamma$ sống
theo Bổ đề \ref{bd:tongkhong} --- đạo hàm của $R$ tại $c$ là
\begin{equation}
DR(c)\big|_H \;=\; \frac{n}{\|c\|_1}\,I .
\label{eq:DR}
\end{equation}
\end{menhde}

\begin{proof}
Hệ quả trực tiếp của Mệnh đề \ref{md:affine}: ánh xạ là affine với phần tuyến
tính $-\tfrac{n\beta}{\|c\|_1}I$ hạn chế trên $H$.
\end{proof}

Một \emph{bội vô hướng của đồng nhất thức}: sai số đầu ra tỉ lệ \emph{đúng}
với $-\Gamma$. Hướng nguyên vẹn; chỉ độ lớn bị nén theo hệ số $n/\|c\|_1$.

\begin{menhde}[Tái lập nhãn từ hướng]
\label{md:tailap}
\textsc{[cm]}\; Giả sử $\bar\varepsilon_j\equiv\bar\varepsilon$ như nhau giữa
các lớp. Từ hướng của $\Gamma$ ta khôi phục được $N$ như sau:
\begin{enumerate}
\item Giải $u=-\bigl(I-Q^{\!\top}\bigr)^{+}\Gamma$ bằng giả nghịch đảo --- cần
  giả nghịch đảo vì $I-Q^{\!\top}$ suy biến (Hệ quả \ref{hq:hatnhan}). Khi đó
  $u=\bar\varepsilon N+t\one$ với $t$ chưa biết.
\item \emph{Nếu có ít nhất một lớp vắng mặt} $j_0$ thì $N_{j_0}=0$ nên
  $u_{j_0}=t$, tức $t=\min_j u_j$, và $N\propto u-t\one$ xác định duy nhất.
\item \emph{Nếu không có lớp vắng mặt}, phép neo $t=\min_j u_j$ trừ đi
  $N_{\min}:=\min_j N_j$, cho
  \begin{equation}
  \hat N_j \;=\; \frac{n\,(N_j-N_{\min})}{\,n-C\,N_{\min}\,}.
  \label{eq:khongvang}
  \end{equation}
\end{enumerate}
\end{menhde}

\begin{proof}
Bước 1 là định nghĩa giả nghịch đảo cộng Hệ quả \ref{hq:hatnhan} (hạt nhân
một chiều, sinh bởi vectơ dừng của $Q$; với $Q$ đối xứng hóa đều thì là
$\one$). Bước 3: $u_j=\bar\varepsilon N_j+t'$ nên $\min_j u_j$ đạt tại
$\arg\min_j N_j$; trừ đi rồi chuẩn hóa về tổng $n$ với
$\sum_j(N_j-N_{\min})=n-CN_{\min}$.
\end{proof}

\begin{nhanxet}[Kết luận của mục này]
\label{nx:nhieu}
Trong \eqref{eq:Gid}, $Q$ mang \emph{hình dạng} của lỗi và $\|a\|_1=n\epstr$
mang \emph{độ lớn}. Mệnh đề \ref{md:tailap} chỉ dùng hướng của $\Gamma$ ---
không dùng $\|\Gamma\|_1$ ở bất kỳ bước nào. Nên:

\begin{quote}
\emph{Độ lớn của tín hiệu là một \textbf{tham số nhiễu}: nó mang thông tin về
$\epstr$, thứ kẻ tấn công không cần. Toàn bộ thông tin nhãn nằm trong hướng.}
\end{quote}

Nhóm phụ trợ mắc lỗi coi độ lớn là tín hiệu: bằng cách so $\Gamma$ với một
hằng số tuyệt đối, chúng để một tham số nhiễu điều khiển đầu ra. Đó là lý do
chúng sụp, và cũng là lý do sụp đổ đó \emph{tránh được}.

Chất lượng khôi phục thực tế của cách đọc theo hướng trên pipeline thật ---
\textbf{chưa đo}.
\end{nhanxet}

% ===========================================================================
\section{Kết luận và nguyên lý thiết kế}
\label{sec:nguyenly}

\subsection{Phát biểu}

\begin{quote}
Huấn luyện làm $\|\Gamma\|_1$ co theo $\epstr$ --- điều này \textbf{bắt buộc},
vì gradient lớp cuối là tổng các phần dư $p_i-e_{y_i}$ nên độ lớn của nó
\emph{chính là} lỗi (Bổ đề \ref{bd:dolon}). Phương pháp dùng phụ trợ so
$\Gamma$ với một thước \emph{cố định không co}: thước của ZLG$^+$ là
$1-\epsaux$ nên tăng về trần, thước của RLU là $\bar\rho\cdot$softmax nên bị
ghim bởi chính phép chuẩn hóa softmax (Mệnh đề \ref{md:ghim}). Tỉ số
$A=\Theta(\epstr)$ co lại, và đầu ra bị \emph{nhốt} trong quả cầu bán kính
$A+C$ quanh tiên nghiệm lớp của tập phụ trợ (Mệnh đề \ref{md:bay}) --- một
ràng buộc đúng ở mọi mức sai số, không phải một ngưỡng --- rồi hội tụ hẳn về
tâm đó khi $A\to0$ (Mệnh đề \ref{md:khongbatbien}). Phương pháp không dùng phụ trợ dựng
thước \emph{từ chính $\Gamma$} --- LLG từ \code{negative\_gradient}, VLI$^*$
từ \code{g\_max} --- nên thước co cùng tín hiệu, tỉ số không đổi, và chúng
bất biến thang độ (Mệnh đề \ref{md:llg}, \ref{md:vlistar}). Thông tin nhãn
nằm ở \emph{hướng} của $\Gamma$ (Mệnh đề \ref{md:huong}, \ref{md:tailap}), thứ
không co.
\end{quote}

Nói ngắn: \textbf{suy giảm không phải vì tín hiệu yếu, mà vì được đo bằng một
cái thước không chịu yếu theo.}

\subsection{Nguyên lý}

\begin{quote}
Điều kiện cần để một tấn công khôi phục nhãn không suy giảm theo huấn luyện
là: đầu ra \textbf{bất biến dưới $\Gamma\mapsto\lambda\Gamma$ với mọi
$\lambda>0$}.
\end{quote}

Đây là điều kiện \emph{kiểm được} bằng một vòng lặp vài dòng trên hàm tấn
công, không cần chứng minh và không cần dữ liệu. Nó là đặc tả đúng của
VLI$^*$: \code{signal/g\_max} là toàn bộ cách VLI$^*$ chạm vào tín hiệu.

\subsection{Tự phản biện: bốn chỗ kết luận này không phủ}

\begin{enumerate}
\item \textbf{Huấn luyện không chỉ làm co $\Gamma$, nó còn đổi hướng.} $Q$
  trôi theo huấn luyện. Bất biến thang độ loại bỏ chế độ hỏng trội, không
  phải mọi chế độ. LLG và VLI$^*$ vẫn có thể suy giảm vì hướng trôi hoặc vì
  $Q$ lệch khỏi giả định đối xứng của VLI$^*$.

\item \textbf{Bất biến giải thích ổn định, không giải thích chất lượng.} Đầu
  ra của LLG bị quyết định bởi $h_1$ và dãy \code{argmin} tham lam, nên nó có
  thể vốn không chính xác cao ngay từ đầu. Quan sát ``LLG không giảm'' hoàn
  toàn tương thích với một phương pháp tầm thường nhưng vững.

\item \textbf{VLI$^*$ đổi chế độ hỏng chứ không xóa.} Nó cần
  \code{g\_max}$>0$, ngầm cần ít nhất một lớp vắng mặt; không có thì
  \eqref{eq:khongvang} cho câu trả lời bị lệch $N_{\min}$.

\item \textbf{Vectơ giới hạn phụ thuộc phụ trợ.} \emph{Việc} hội tụ về một
  vectơ không phụ thuộc nhãn là chứng minh được; \emph{vectơ nào} thì do phụ
  trợ quyết định ($n\bar p/\sum\bar p$ cho ZLG$^+$,
  $n\cdot\mathrm{softmax}(\bar\mu)$ cho RLU).

\item \textbf{Chặn bẫy chỉ cắn khi client lệch khỏi tiên nghiệm phụ trợ.}
  Hệ quả \ref{hq:chanduoi} chặn dưới sai số bằng
  $\|N-\quant(R(c))\|_1-(A+C)$. Nếu client tình cờ có phân bố lớp gần với
  tiên nghiệm phụ trợ thì số hạng đầu nhỏ và chặn không nói gì --- ZLG$^+$,
  RLU sẽ \emph{trông như} vẫn chạy tốt dù không suy ra được gì từ nhãn. Đây
  là một tiên đoán phân biệt được: hiệu năng còn lại của nhóm phụ trợ phải
  tương quan với \emph{độ lệch của client khỏi phân bố đều}, chứ không với
  chất lượng suy luận. Với Dirichlet $\alpha=0{,}5$ thì $N$ rất lệch nên chặn
  cắn mạnh, nhưng điều đó cần đo chứ không được giả định.

\item \textbf{Chặn bẫy không đóng được hồ sơ RLU ở sai số cao.} Theo Nhận xét
  \ref{nx:rluyeu}, $\bar\rho\approx1{,}1$ làm chặn của RLU rỗng khi
  $\epstr\gtrsim0{,}55$. Kết luận ``nhóm phụ trợ suy giảm sớm'' hiện được
  chứng minh \emph{chặt cho ZLG$^+$} nhưng chỉ \emph{tiệm cận cho RLU}. Đây
  là chỗ yếu nhất còn lại, và nó chỉ đóng được bằng số đo $\|\Gamma\|_1$
  thật, không đóng được bằng thêm đại số.
\end{enumerate}

\subsection{Những gì cần đo}

Danh sách các đại lượng báo cáo này ghi là \textbf{chưa đo}, theo thứ tự ưu
tiên:

\begin{enumerate}
\item \textbf{Kiểm ràng buộc cứng \eqref{eq:bay}.} Đo $A$ và
  $\|\hat N-\quant(R(c))\|_1$ theo từng cửa sổ. Bất đẳng thức
  $\|\hat N-\quant(R(c))\|_1\le A+C$ phải đúng ở \emph{mọi} cửa sổ; một vi
  phạm duy nhất bác bỏ dạng chuẩn \eqref{eq:dangchuan}. Đây là phép thử mạnh
  nhất trong danh sách vì nó là ràng buộc nhị phân, không phải tương quan.
  Kèm theo: $A$ so với $\epstr$ (tiên đoán: tỉ lệ thuận), và irec so với $A$
  --- nếu irec tụt \emph{nhanh hơn} $A$ thì có cơ chế thứ hai chưa nhận ra.

\item \textbf{$\|\Gamma\|_1$ thật so với chặn $2n\epstr$.} Quyết định liệu
  Mệnh đề \ref{md:bay} có cắn được cho RLU ở $\epstr$ cao hay không (Nhận xét
  \ref{nx:rluyeu}). Nếu $\|\Gamma\|_1/(2n\epstr)$ nhỏ hơn $\bar\rho/2$ thì hồ
  sơ RLU đóng lại; nếu không, cần một cơ chế khác.

\item \textbf{Tiên đoán $1/C$} (Nhận xét \ref{nx:lientuc}): chạy cùng
  $\epstr$ trên CIFAR10 và CIFAR100, kiểm $A/n$ của ZLG$^+$ có giảm khoảng
  một bậc không. RLU \emph{không} được giảm theo cách đó --- nếu cả hai cùng
  giảm như nhau thì $1/C$ không phải cơ chế thật.

\item \textbf{Độ lệch của client} (tự phản biện 5): tương quan giữa irec còn
  lại của ZLG$^+$/RLU và $\|N-\quant(R(c))\|_1$. Nếu tương quan âm mạnh thì
  Hệ quả \ref{hq:chanduoi} là lời giải thích đúng; nếu không, chặn bẫy không
  phải cơ chế trội.
\item $w_{\max}=\max_i s_i/\bar O$ --- giả thiết duy nhất còn lại của Bổ đề
  \ref{bd:dolon} cho $\tilde G$.
\item Hành vi của $\bar O$ khi huấn luyện (Nhận xét \ref{nx:Obar}).
\item $\min_j\mathrm{softmax}(\bar\mu)_j$ --- chỉ quyết định dạng affine
  \emph{chính xác} của RLU có dùng được không (Hệ quả \ref{hq:nguong}); không
  ảnh hưởng chặn bẫy (Nhận xét \ref{nx:vaitroH2}), nên ưu tiên thấp.
\item $\epstr/\epsaux$ --- đại lượng điều khiển của LLG$^+$ (Nhận xét
  \ref{nx:llgpratio}).
\end{enumerate}

\appendix
\section{Vị trí trong mã nguồn của mọi khẳng định \textsc{[cm]}}

\begin{center}
\begin{tabular}{@{}lll@{}}
\toprule
\textbf{Khẳng định} & \textbf{Tệp} & \textbf{Hàm} \\
\midrule
ZLG$^+$ có dạng \eqref{eq:dangchuan} & \code{client/attacks.py} & \code{ZLGp\_attack} \\
RLU có dạng \eqref{eq:dangchuan} & \code{client/client\_utils.py} & \code{estimated\_entropy\_from\_grad} \\
$c_{\text{RLU}}=\bar\rho\,\mathrm{softmax}(\bar\mu)$ & \code{client/attacks.py} & \code{RLU\_attack} \\
$\bar p_j=1-\bar\varepsilon_{\mathrm{aux},j}$ & \code{client/client\_utils.py} & \code{estimate\_static\_ZLG} \\
$\bar O=\mathrm{mean}(s_i)$ trên phụ trợ & \code{client/client\_utils.py} & \code{estimate\_static\_ZLG} \\
LLG bất biến & \code{client/attacks.py} & \code{LLG\_attack} \\
VLI$^*$ bất biến & \code{client/attacks.py} & \code{VLI\_star\_closed\_form} \\
LLG$^+$ không bất biến & \code{client/attacks.py} & \code{LLGp\_attack} \\
offset của LLG$^+$ là gradient phụ trợ & \code{client/client\_utils.py} & \code{estimate\_static\_LLG} \\
$\bar\rho$ không phụ thuộc mô hình & \code{client/attacks.py} & \code{rlu\_rho\_mean} \\
Phụ trợ lấy từ test set & \code{main.py} & dòng \code{aux\_dataset = ...} \\
\bottomrule
\end{tabular}
\end{center}

\end{document}
